\documentclass[11pt]{article}
\usepackage{amsmath, amssymb, amsthm}
\usepackage{geometry}
\usepackage{hyperref}
\usepackage{xcolor}
\usepackage{listings}
\geometry{margin=1in}

\newtheorem{proposition}{Proposition}
\newtheorem{remark}{Remark}

\newcommand{\HN}{\mathrm{HN}}
\newcommand{\N}{\mathcal{N}}
\newcommand{\Phi}{\Phi}
\newcommand{\phione}{\phi_1}
\newcommand{\phitwo}{\phi_2}
\newcommand{\zstar}{z^*}
\newcommand{\ztstar}{z^*_t}
\newcommand{\zinit}{z^{(0)}}
\newcommand{\sig}{\sigma}
\newcommand{\mh}{\delta}

\title{Root Cause and Fix: Frozen $\phi_1^z$, $\phi_2^z$ in the\\
       AR(2) Temporal Skew-$t$ Sampler}
\author{}
\date{2026-05-27}

\begin{document}
\maketitle

\section{Background: The Half-Normal Copula Parameterisation}

The spatial skew-$t$ model introduces a latent scale process $z_{k,t} > 0$
for knot $k \in \{1,\ldots,K\}$ and time $t \in \{1,\ldots,T\}$.
Conditionally on the precision $\tau_{k,t}$, the marginal prior is

\begin{equation}
  z_{k,t} \mid \tau_{k,t} \;\sim\; \HN\!\left(\sigma_{k,t}\right),
  \qquad \sigma_{k,t} = \tau_{k,t}^{-1/2},
  \label{eq:z_prior}
\end{equation}

where $\HN(\sigma)$ denotes the half-normal distribution with scale $\sigma$,
i.e.\ the distribution of $|\epsilon|$ for $\epsilon \sim \N(0,\sigma^2)$,
with CDF
\[
  F_{\HN}(z;\,\sigma) \;=\; 2\,\Phi\!\left(\tfrac{z}{\sigma}\right) - 1,
  \qquad z \ge 0.
\]

To impose AR(2) temporal dependence while respecting the positivity
constraint, the sampler works in the copula-transformed space.
Define the \emph{half-normal copula map} $h_\sigma : [0,\infty) \to \mathbb{R}$,
\begin{equation}
  \zstar_{k,t} \;=\; h_{\sigma_{k,t}}(z_{k,t})
  \;:=\; \Phi^{-1}\!\bigl(F_{\HN}(z_{k,t};\,\sigma_{k,t})\bigr)
  \;=\; \Phi^{-1}\!\!\left(2\,\Phi\!\left(\tfrac{z_{k,t}}{\sigma_{k,t}}\right)-1\right),
  \label{eq:copula}
\end{equation}
with inverse
\[
  z_{k,t} \;=\; h_{\sigma_{k,t}}^{-1}(\zstar_{k,t})
  \;=\; \sigma_{k,t}\,\Phi^{-1}\!\!\left(\tfrac{1+\Phi(\zstar_{k,t})}{2}\right).
\]
By construction, $\zstar_{k,t} \in \mathbb{R}$ and, under the marginal prior
\eqref{eq:z_prior}, $\zstar_{k,t} \sim \N(0,1)$.

The temporal prior on $\zstar_{k,\cdot}$ is a stationary AR(2):
\begin{equation}
  \zstar_{k,t} \;=\; \phione\,\zstar_{k,t-1} \;+\; \phitwo\,\zstar_{k,t-2}
  \;+\; \varepsilon_{k,t},
  \qquad
  \varepsilon_{k,t} \stackrel{\mathrm{iid}}{\sim} \N(0,\,\sigma_\varepsilon^2),
  \label{eq:ar2}
\end{equation}
where, normalising so that $\gamma_0 := \mathrm{Var}(\zstar_{k,t}) = 1$,
the Yule--Walker equations give
\[
  \gamma_1 = \frac{\phione}{1-\phitwo},
  \quad
  \gamma_2 = \phione\gamma_1 + \phitwo,
  \quad
  \sigma_\varepsilon^2 = 1 - \phione\gamma_1 - \phitwo\gamma_2.
\]

\section{The MCMC Update Structure}

At each iteration the sampler executes two linked Metropolis--Hastings
steps (implemented in \texttt{updateZTS\_AR2} and \texttt{updatePhiAR2TS}).

\subsection*{Step 1 — Update $z_{k,t}$ (equivalently $\zstar_{k,t}$)}

For each $(k,t)$, propose
\[
  \zstar_{k,t}^{\mathrm{can}} \;\sim\; \N\!\left(\zstar_{k,t},\,\mh_{k,t}^2\right),
\]
transform back via $z_{k,t}^{\mathrm{can}} = h_{\sigma_{k,t}}^{-1}(\zstar_{k,t}^{\mathrm{can}})$,
and accept with log-ratio
\begin{align}
  R_z \;=\;&
  \log p\!\left(y_{\cdot,t} \mid z^{\mathrm{can}}_{k,t},\,\ldots\right)
  -\log p\!\left(y_{\cdot,t} \mid z_{k,t},\,\ldots\right) \notag\\
  &+\;\log f_{\AR}\!\left(\zstar_{k,t}^{\mathrm{can}} \mid \zstar_{k,t-1},\zstar_{k,t-2};\,\phione,\phitwo\right)
  -\log f_{\AR}\!\left(\zstar_{k,t} \mid \zstar_{k,t-1},\zstar_{k,t-2};\,\phione,\phitwo\right) \notag\\
  &+\;\text{(forward-looking terms for $t+1$ and $t+2$)},
  \label{eq:Rz}
\end{align}
where $f_{\AR}(\cdot\mid\cdot)$ is the conditional density of \eqref{eq:ar2}.

\subsection*{Step 2 — Update $(\phione,\phitwo)$}

Given the current $\zstar_{k,\cdot}$, the transition log-likelihood is
\begin{equation}
  \ell(\phione,\phitwo) \;=\;
  \sum_{k=1}^{K}\sum_{t=3}^{T}
  \log \N\!\left(\zstar_{k,t};\;\phione\,\zstar_{k,t-1}+\phitwo\,\zstar_{k,t-2},\;
  \sigma_\varepsilon^2(\phione,\phitwo)\right).
  \label{eq:phi_ll}
\end{equation}
A Metropolis step proposes $\phione^{\mathrm{can}} \sim \N(\phione,\mh_{\phi_1}^2)$
and accepts if the candidate lies in the AR(2) stationarity region
$\mathcal{S} = \{|\phitwo|<1,\;\phione+\phitwo<1,\;\phitwo-\phione<1\}$
and the log-ratio
\[
  R_\phi \;=\;
  \ell(\phione^{\mathrm{can}},\phitwo) - \ell(\phione,\phitwo)
  +\log\pi(\phione^{\mathrm{can}}) - \log\pi(\phione)
\]
satisfies $\log U < R_\phi$ for $U \sim \mathrm{Uniform}(0,1)$.

\section{The Failure Mode: $\zinit = 0$}

\subsection{Singularity of the Copula Map at the Boundary}

The copula map $h_\sigma$ has a \emph{non-integrable singularity} at $z = 0$:
\begin{equation}
  \lim_{z\to 0^+} h_\sigma(z)
  \;=\; \Phi^{-1}\!\bigl(2\,\Phi(0) - 1\bigr)
  \;=\; \Phi^{-1}(0)
  \;=\; -\infty.
  \label{eq:singularity}
\end{equation}
Consequently, initialising $z_{k,t}^{(0)} = 0$ for all $k,t$ yields
\[
  \zstar_{k,t}^{(0)} \;=\; h_{\sigma_{k,t}}(0) \;=\; -\infty
  \qquad \forall\, k,\, t.
\]

\subsection{Cascading Failure in Step 1}

With $\zstar_{k,t}^{(0)} = -\infty$, the proposal for $z_{k,t}$ at the
first iteration draws
\[
  \zstar_{k,t}^{\mathrm{can}} \;\sim\; \N(-\infty,\,\mh_{k,t}^2) \equiv \mathrm{NaN}.
\]
Because $\mathrm{NaN}$ propagates through $h_\sigma^{-1}$ and all
downstream arithmetic, the log-ratio \eqref{eq:Rz} evaluates to
$R_z = \mathrm{NaN}$.  The guard
\[
  \texttt{if (!is.na}(R_z)\texttt{)}\;\{\;\text{accept}\;\}
\]
is therefore \emph{never triggered}, and $z_{k,t}$ remains at $0$
\textbf{for every subsequent iteration}.

\begin{proposition}
  If $z_{k,t}^{(0)} = 0$ for all $k,t$, then
  $z_{k,t}^{(m)} = 0$ for all iterations $m \ge 1$
  under the update rule of \texttt{updateZTS\_AR2}.
\end{proposition}

\subsection{Propagation to the $\phi$ Update (Step 2)}

Since $z_{k,t}^{(m)} = 0$ for all $m$, the copula-transformed chain
satisfies $\zstar_{k,t}^{(m)} = -\infty$ for all $m$.
The transition log-likelihood \eqref{eq:phi_ll} at the current state
therefore evaluates to
\[
  \ell(\phione,\phitwo) \;=\;
  \sum_{k,t} \log\N(-\infty;\;\phione\cdot(-\infty)+\phitwo\cdot(-\infty),\;\sigma_\varepsilon^2)
  \;=\; -\infty.
\]
For any candidate $\phione^{\mathrm{can}}$ in the stationarity region
$\mathcal{S}$, the candidate log-likelihood likewise involves terms of
the form $-\infty - (-\infty)$, which is indeterminate:
\[
  \ell(\phione^{\mathrm{can}},\phitwo) - \ell(\phione,\phitwo)
  \;=\; (-\infty) - (-\infty)
  \;=\; \mathrm{NaN}.
\]
Hence $R_\phi = \mathrm{NaN}$ for every proposal, and the acceptance
guard is never satisfied.

\begin{proposition}
  If $z_{k,t}^{(0)} = 0$ for all $k,t$, then $\phione^{(m)} = \phione^{(0)} = 0$
  and $\phitwo^{(m)} = \phitwo^{(0)} = 0$ for all $m \ge 1$.  The AR(2)
  coefficient chain is permanently frozen at its initial value.
\end{proposition}

\begin{remark}
  The failure is specific to the copula-space updater.  The $\tau$
  update (\texttt{updateTauTS\_AR2}) uses a gamma copula whose argument
  is $\tau^{(0)} = 1 > 0$, giving a finite $\tau^*$.  The knot update
  (\texttt{updateKnotsTS\_AR2}) uses a probit transform of spatial
  coordinates, which are also bounded away from the boundary.  Neither
  is affected.
\end{remark}

\section{The Fix: Principled Initialisation at the Marginal Median}

\subsection{Choice of $\zinit$}

Any $\zinit > 0$ breaks the singularity \eqref{eq:singularity}.  The
most natural choice is the \emph{median} of the marginal prior
\eqref{eq:z_prior}, which satisfies
\[
  F_{\HN}\!\left(\zinit;\,\sigma\right) = \tfrac{1}{2}
  \;\;\Longleftrightarrow\;\;
  2\,\Phi\!\left(\tfrac{\zinit}{\sigma}\right)-1 = \tfrac{1}{2}
  \;\;\Longleftrightarrow\;\;
  \zinit = \sigma\,\Phi^{-1}\!\!\left(\tfrac{3}{4}\right).
\]
Writing $q_{3/4} := \Phi^{-1}(3/4) \approx 0.6745$, this gives
\begin{equation}
  \zinit \;=\; \frac{q_{3/4}}{\sqrt{\tau_{\mathrm{init}}}},
  \qquad q_{3/4} \approx 0.6745.
  \label{eq:z_init_fix}
\end{equation}

\subsection{Why the Median is the Right Anchor}

The copula map evaluated at $\zinit$ from \eqref{eq:z_init_fix} yields

\begin{equation}
  h_\sigma(\zinit)
  \;=\; \Phi^{-1}\!\bigl(F_{\HN}(\zinit;\,\sigma)\bigr)
  \;=\; \Phi^{-1}\!\!\left(\tfrac{1}{2}\right)
  \;=\; 0,
  \label{eq:zstar_at_median}
\end{equation}

\emph{regardless of $\sigma$ (equivalently, of $\tau_{\mathrm{init}}$)}.
The chain therefore starts at the centre of the $\zstar$ space,
minimising asymmetry in the initial MH proposals for both $z$ and $\phi$.

\subsection{Scale Invariance}

The result \eqref{eq:zstar_at_median} holds for every value of
$\tau_{\mathrm{init}}$:

\begin{center}
\begin{tabular}{cccc}
\hline
$\tau_{\mathrm{init}}$ & $\sigma = \tau_{\mathrm{init}}^{-1/2}$ &
$\zinit = q_{3/4}\,\sigma$ & $\zstar^{(0)} = h_\sigma(\zinit)$ \\
\hline
0.1  & 3.162 & 2.133 & 0 \\
1    & 1.000 & 0.6745 & 0 \\
4    & 0.500 & 0.337  & 0 \\
16   & 0.250 & 0.168  & 0 \\
\hline
\end{tabular}
\end{center}

By contrast, the previous default $\zinit = 1$ gives
$\zstar^{(0)} = h_\sigma(1) \approx 0.47$ when $\tau_{\mathrm{init}} = 1$,
and $\zstar^{(0)} \approx 1.69$ when $\tau_{\mathrm{init}} = 4$ —
increasingly far from the centre as the scale deviates from 1.

\subsection{Implementation}

In \texttt{code/R/ar2/mcmc\_ar2.R}, the function signature changes from
\[
  \texttt{z.init = 0} \;\longrightarrow\; \texttt{z.init = NULL},
\]
and the initialisation block computes
\begin{lstlisting}[language=R, basicstyle=\ttfamily\small]
if (is.null(z.init)) {
    z.init <- 0.6745 / sqrt(tau.init)   # median of HN(1/sqrt(tau.init))
}
z <- matrix(z.init, nknots, nt)
\end{lstlisting}
Callers who supply \texttt{z.init} explicitly are unaffected.

\section{Observed Symptom in the Simulation Study}

In simulation settings 9 and 10 (skew-$t$, $K=1$ and $K=5$,
AR(2) temporal; analysis methods 7 and 8), the posterior summary CSV files
showed
\[
  \hat{\phi}_1^z = \hat{\phi}_2^z = 0,
  \qquad
  \widehat{\mathrm{SD}}_{\mathrm{post}}(\phi_1^z)
  = \widehat{\mathrm{SD}}_{\mathrm{post}}(\phi_2^z) = 0,
\]
across all 10 datasets, confirming that every post-burn draw was exactly
equal to the initialisation value.  Concurrently,
$\hat{\phi}_1^\tau, \hat{\phi}_2^\tau$ were correctly estimated
(settings 9: $\hat\phi_1^\tau \approx 0.88$, truth $\phi_1 = 0.80$),
consistent with the diagnosis in Section~3.

\end{document}
